\documentclass[11pt,a4paper]{article}
\usepackage[T1]{fontenc}
\usepackage{lmodern,amsmath,amssymb,booktabs,microtype}
\usepackage[margin=24mm]{geometry}
\usepackage[colorlinks=true,urlcolor=blue,linkcolor=blue]{hyperref}
\newcommand{\db}{\partial_{\overline H}}
\newcommand{\CC}{\mathbb C}
\title{Literal weak integrals and\\moving-pole continuation}
\author{Alan Day \and Chengheng\thanks{AI research assistant (OpenAI).}}
\date{Working research note --- 1 November 2026}
\begin{document}\maketitle
\begin{abstract}
A locally integrable single-chiral multiplier need not define a
holomorphic family of weak pairings. Applying the standard distributional
Cauchy-kernel identity to the known quantum-dilogarithm pole divisor,
we compute the parameter-space source term exactly. A compact local
example separates the literal weak integral from the meromorphic
continuation of its exterior values, and a fixed normalized test state
has the same obstruction. Full-product quadrature checks the formulas.
The result concerns one specified weak prescription in the interior
of the first quadrant. It neither refutes the final bounded pentagon
nor constructs a physical de Sitter observer state.
\end{abstract}
\section{Existing ingredients and the question}
Use dimensionless real coordinates $(x,y)$, Lebesgue measure, and
$H=a+ib$ with $a,b>0$. The known scalar function \cite{ks} is
\[
 \Phi^H(z)=\frac{\psi_q(e^z)}{\psi_{\widetilde q}(e^{z/H})},
 \quad q=e^{\pi iH},\quad\widetilde q=e^{-\pi i/H},
 \quad \psi_q(v)=\prod_{k\geq0}(1+q^{2k+1}v)^{-1}.
\]
Its first simple pole is $z_0(H)=-\pi i(1+H)$, with residue
\[
 R_H=-\frac{(\widetilde q^2;\widetilde q^2)_\infty}
             {(q^2;q^2)_\infty}\ne0,\qquad
 (v;q)_\infty=\prod_{k\geq0}(1-vq^k).
\]
The residue is holomorphic and nonzero on the upper half-plane.
We ask whether the literal compact pairing
\[
 W_f(H)=\int_{\mathbb R^2}f(x,y)\Phi^H(x+Hy)\,dx\,dy
\]
is holomorphic when the pole moves through the support of a
\emph{fixed} test density $f$.

The main mathematical input is standard:
$\partial_{\bar z}(1/z)=\pi\delta_0$ for
$\partial_{\bar z}=(\partial_{\Re z}+i\partial_{\Im z})/2$ and
ordinary area measure \cite{ramond}. Thus pointwise holomorphy
away from a moving pole does not justify differentiating an integral
as though that pole were absent. We use this established identity,
not a newly proposed continuation theorem.

\newpage
\section{A precise local pairing and its source term}
Let
\begin{align}
 p(H)&=(x_0(H),y_0(H))
 =\left(\frac{\pi(a^2+a+b^2)}b,-\frac{\pi(a+1)}b\right),\\
 u(x,y)&=-\frac{x+\pi i}{y+\pi i},\qquad
 w(x,y)=y+\pi i,\\
 d_H(x,y)&=x+Hy+\pi i(1+H)=w(x,y)(H-u(x,y)).
\end{align}
The pole map is a diffeomorphism onto
$\{x>0,y<0,xy<-\pi^2\}$; its positive real Jacobian is
\[
 J(H)=\det Dp(H)=\frac{\pi^2((a+1)^2+b^2)}{b^3}
             =\frac{|w(p(H))|^2}{b}.
\]
Choose small neighborhoods $V$ in parameter space and $U$ in coordinate
space such that only this pole branch meets $V\times U$.
Then
\[
 \Phi^H(x+Hy)=\frac{R_H}{d_H(x,y)}+E(H,x,y),
\]
where $E$ is holomorphic in $H$, locally uniformly for $(x,y)$ in
compact subsets of $U$. For $f\in C_c^\infty(U)$ the singularity is
locally integrable in two real variables. Applying the Cauchy identity
and changing variables through $p$ gives
\begin{align}
 \db W_f(H)
 &=\pi R_H\frac{J(H)}{w(p(H))}\,f(p(H))\\
 &=\frac{\pi R_H}{b}
       \bigl(y_0(H)-\pi i\bigr)f(p(H)).
 \label{source}
\end{align}
All equalities hold distributionally in $H$. One justification is
Fubini on the locally integrable joint kernel: at fixed $(x,y)$,
\[
 \db\frac{1}{w(H-u)}=\frac{\pi}{w}\delta^{(2)}(H-u).
\]
The holomorphic remainder contributes zero. The conjugate factor in
\eqref{source} is essential; the Jacobian cannot be omitted.
Whenever $f(p(H))\ne0$, the source is nonzero.

\subsection*{An explicit region with no other poles}
Take $H_c=1+i$ and $V=\{|H-H_c|<.08\}$. Put the source support
inside $p(\{|u-H_c|\leq.04\})$. Both lie over the square
$Q=[.9,1.1]^2$ in $(a,b)$.
The other poles have coordinates
\[
 y_{mn}(H)=-\frac{\pi((2m+1)+a(2n+1))}{b},\quad m,n\geq0 .
\]
First-pole support has
$-7\pi/3\leq y\leq-19\pi/11$.
If $m+n\geq1$, then $y_{mn}\leq-37\pi/11<-7\pi/3$.
Thus no other pole enters this compact test. This isolation is exact,
not inferred from a finite scan of poles.

\newpage
\section{Literal values and exterior continuation differ}
Let $r_0=.04$, $z=H-H_c$, and
\[
 g(u)=\left(1-\frac{|u-H_c|^2}{r_0^2}\right)_+^8,
 \qquad m_g=\int g\,dA=\frac{\pi r_0^2}{9}.
\]
This compact function is $C^7$; the locally integrable calculation above
also applies to it. Choose the fixed complex density
\[
 f(p(u))=\frac{w(p(u))}{J(u)}g(u).
\]
It is a test pairing, not a probability density. Changing variables
exactly, the residue-normalized pairing becomes
\[
 \mathcal W(H):=\frac{W_f(H)}{R_H}
       =\int_{\CC}\frac{g(u)}{H-u}\,dA(u)+K(H)
       =\mathcal C_g(H)+K(H),
\]
where $K$ is holomorphic throughout $V$. Angular integration gives
\[
 \mathcal C_g(H)=
 \begin{cases}
 \displaystyle\frac{\pi r_0^2}{9z}
       \left[1-\left(1-\frac{|z|^2}{r_0^2}\right)^9\right],
       &0<|z|<r_0,\\[6pt]
 0,&z=0,\\[2pt]
 \displaystyle\frac{m_g}{z},&|z|\geq r_0 .
 \end{cases}
\]
In particular $\db\mathcal W=\pi g$, and at $H_c$ this is $\pi$.
The exterior values have the unique meromorphic continuation
$\mathcal W_{\rm ext}=K+m_g/z$ to $V$, with a pole at $H_c$.
The literal value there is finite. For $0<|z|<r_0$ their exact difference is
\[
 \mathcal W-\mathcal W_{\rm ext}
       =-\frac{m_g}{z}
         \left(1-\frac{|z|^2}{r_0^2}\right)^9.
\]
No function holomorphic on all of $V$ can match the exterior values:
on the counterclockwise circle $|z|=.06$,
\[
 \oint\mathcal W(H)\,dH=2\pi i\,m_g\ne0 .
\]
Subtracting the full Cauchy transform gives the holomorphic remainder
$K$, but changes the exterior values too. Adding a holomorphic function
cannot remove a nonzero $\bar\partial$ source.

\subsection*{A normalized nonnegative test state}
Define one fixed real state by
\[
 |\eta(p(u))|^2=\frac{g(u)}{J(u)m_g},\qquad\|\eta\|_{L^2}=1.
\]
It is compact and $C^3$. Its spectral first moment of $\Phi^H(X+HY)$
is finite, while the vector is outside that operator's domain when
the pole meets its nonzero region. For its residue-normalized quadratic
form $\mathcal W_\eta$,
\[
 \db\mathcal W_\eta(H_c)
 =\frac{\pi}{m_g(-2\pi+i\pi)}
 =\frac{-2-i}{5m_g}\ne0.
\]
This is a well-defined first-moment form, not an $L^2$ operator action
on $\eta$, and not a claim that $\eta$ is a physical de Sitter state.

\newpage
\section{Full-product calculations}
We computed the complete $\Phi^H$, rather than only its pole.
For $d=w(p(u))(H-u)$ define
\[
 E_H(d)=
 \frac{(\widetilde q^2e^{d/H};\widetilde q^2)_\infty}
      {(\widetilde q^2;\widetilde q^2)_\infty}
 \frac{(q^2;q^2)_\infty}{(q^2e^d;q^2)_\infty}.
\]
Then $\Phi^H(x+Hy)/R_H=E_H(d)/(e^d-1)$, and the regular integrand is
\[
 w(p(u))\left[\frac{E_H(d)}{\operatorname{expm1}(d)}-\frac1d\right].
\]
We integrate this smooth difference on the fixed source disk and add
the exact $\mathcal C_g$. Independently, for interior $H$ we integrate
the full product in polar coordinates centered at its pole; the area
Jacobian removes the integrable $1/r$ singularity. Exterior checks
use the full product on the fixed disk.

At 40-digit working precision, six points were evaluated using
8/32 and 16/64 radial/angular nodes for subtraction, and 16/64 for
the direct route. The largest route difference was $1.43\,10^{-22}$;
the largest subtraction-grid change was $1.76\,10^{-31}$.
These are empirical comparisons, not certified error bounds.

With four-point central differences for $\partial_{\bar H}$,
the complex test approaches the exact value $\pi$:
\begin{center}
\begin{tabular}{rr}
\toprule
step $\delta$ & relative deviation from $\pi$\\
\midrule
.004 & .03908052789\\
.002 & .00994188494\\
.001 & .00249635760\\
.0005& .00062477219\\
\bottomrule
\end{tabular}
\end{center}
The ratio of successive deviations approaches four, consistent with
second-order differencing. The two smallest steps were also evaluated
by the direct route, agreeing within $1.57\,10^{-34}$ in the derivatives.

The full-product contour integral on $|H-H_c|=.06$ approaches
$2\pi i\,m_g$. For the loop, radial, and source angular orders
$(8,12,32)$, $(16,16,48)$, $(32,16,64)$, relative deviations are
$3.44\,10^{-10}$, $9.14\,10^{-17}$, and $1.54\,10^{-20}$.
For the normalized positive state the derivative approaches
\[
-716.1972439135-358.0986219568i.
\]
Its final relative deviation is $.00062478468$.

There are 13 main records and three positive-state difference records,
with samples, parameters, versions, and input hashes.
All 34 symbolic/numerical diagnostics and 62 saved-arithmetic/hash
replay checks pass. Replay does not re-integrate the quantum dilogarithm;
the analytic Cauchy-kernel argument establishes the non-holomorphy.

\newpage
\section{What this resolves, and what remains}
The local example answers a specific question left by the preceding
domain study: replacing a divergent operator action by its
locally integrable weak pairing does \emph{not}, by itself, justify
fixed-vector holomorphic continuation. The missing contribution is
the ordinary distributional Cauchy source from the moving pole.
This is an application of known complex analysis, not a new general
theorem about quantum dilogarithms.

The statement has several boundaries.
The example lies inside the first parameter quadrant, around $1+i$;
it is not a direct computation of the boundary $H=i\hbar$.
The compact test density is fixed as $H$ varies.
A chosen contour, a parameter-dependent test space, or a
meromorphic continuation can define a different object.
The example proves that such an object cannot silently be identified
with the literal weak integral on the whole neighborhood.

In particular, nothing here computes the bounded paired word
\[
 F_\hbar(X,Y)F_\hbar(P_y,P_x)
 \quad\text{versus}\quad
 B M^{-1} B M A.
\]
The bounded phase remains well-defined almost everywhere even when its
individual chiral factors have these domain and continuation problems.
The original pentagon reproduction is still unresolved.

\paragraph{Next research action.}
Do not continue trying to rescue the single-factor argument by merely
declaring a weak integral. Return to identities or Fourier kernels
for the \emph{paired bounded function} with all measure and representation
hypotheses stated. An existing theorem with those hypotheses would be
preferable to a new proof. Otherwise a singularity-aware certified
bounded matrix element could distinguish the candidate words.
The previous finite-grid residuals are motivation, not certification.

The physical goal remains the same: determine the actual observer
insertion, reflection, mass labels, and gluing measure in a common
Hilbert-space prescription before claiming a positive four-point
Gram matrix or a finite-$G$ holographic dictionary. A local
mathematical obstruction is not a no-go theorem for de Sitter holography.

\begin{thebibliography}{9}\small
\bibitem{ks} H. K. Kim and C. Scarinci,
\emph{A quantization of moduli spaces of 3-dimensional gravity},
\href{https://arxiv.org/abs/2112.13329v2}{arXiv:2112.13329v2},
especially \S6.3, Prop.~6.6 (pole divisor and compact-product formula),
and \S6.4 (the continuation step under examination).
\bibitem{ramond} T. Ramond,
\emph{Distributions and Partial Differential Equations},
lecture notes, version 2.06, 19 January 2015,
Example 4.4.5 and Proposition 4.4.6, printed p.~87 (PDF p.~88).
\href{https://www.math.u-psud.fr/~ramond/docs/cours/distributionsedp.pdf}
{Author's university-hosted notes}.
\end{thebibliography}
\end{document}
